% 자동 생성 — scripts/build_tables.py. 손으로 고치지 않는다.
% 출처: output/results/21_b_group.json : 대칭
\begin{table}[htbp]
  \centering
  \caption{두 사건을 같은 자로. 같은 통제군 정의, 같은 창, 같은 추정량.}
  \label{tab:bgroup}
  \begin{tabular}{lrrrr}
    \toprule
    \textbf{사양} & \textbf{처치 동} & \textbf{효과 (\%)} & \textbf{WCR $p$} & \textbf{무작위화 $p$} \\
    \midrule
    A 2020.6 (처치 4개 동) & 4 & -49.2 & 0.038 & $<0.001$ \\
    B 2021.4 (처치 5개 동) & 5 & -44.7 & 0.057 & $<0.001$ \\
    B 2021.4 (목동·여의도만) & 2 & -46.2 & 0.136 & 0.001 \\
    \bottomrule
  \end{tabular}

  \begin{minipage}{\linewidth}\footnotesize
  \textit{주.} 창 전체($\pm$12개월) 이중차분이다. 늘 인용하는 $-79.8\%$ 는 지정 직후 3개월이라 재는 것이 다르다. 두 사건의 차이를 쌓아서 직접 추정하면 $+8.7\%$, 무작위화 $p=0.715$ 다.
  \end{minipage}
\end{table}
